\scanpage{147}\mspage{2}
\={m} des $2$-cat. tout parait si clair -- ``il n'y a qu'à suivre le bout de son nez''\dots) et enfin l'absence d'une définition d'un foncteur allant des ens. (semi) simpliciaux vers les $n$-groupoïdes. Ce foncteur correspond à une intuition géométrique si claire qu'une théorie qui ne l'inclurait pas me semblerait canularesque. Peut-être en \inserted{essayant d'}écri\deleted{\unsure{vant}}\unsure{re} (Comme\unsure{.} une sorte de liste de voeux de Noël !) de façon assez complète et explicite les \deleted{relations} notions \deleted{\unsure{qu'on}}dont on voudrait disposer et les relations (foncteurs, équivalences etc) qui doivent les lier, arriverait-on finalement à une sorte de description axiomatique assez poussée \deleted{\unsure{pour}}qui, soit \unsure{donne|donner} la clef d'une construction ad-hoc explicite, soit \inserted{qui permette} à la rigueur d'énoncer et démontrer un th. d'existence et d'unicité d'une théorie du type souhaité.

Par ailleurs, n'ayant pas compris l'idée de Segal dans ta dernière lettre (que j'ai généreusement envoyée à Illusie\dots), je ne vois pas comment tu définis les $n$-cat. de Picard -- mais peu importe. Pour ce qui est des $n$-cat. de Picard ``strictes'', tout ce que je leur demande, c'est qu'ils forment finalement une \inserted{$(n+1)$-}catégorie $(n+1)$-équivalente à celle des complexes de chaines de longueur $n$. D'accord ? Je te remercie
% archive sheet number: 147
